\section{四层零依赖工程架构与技术选型考量}\label{sec:architecture}

在动手搭建前端之前，我们小组认真讨论过一个很现实的问题：到底要不要上 React + Vite + Node 后端服务？讨论之后我们一致投了反对票。原因很简单：课程验收时老师和助教需要在不同电脑上快速拉起项目，如果强依赖特定版本的 Node.js、复杂的 \texttt{npm install} 编译链甚至在线 CDN 脚本，一旦遇到校园网波动或环境版本冲突，整个演示就会卡在安装依赖上。

基于这一工程判断，我们将全系统划分为职责边界清晰的\textbf{四层解耦架构}（图~\ref{fig:architecture}），实现了“离线重计算在 Python、在线轻交互在原生浏览器”的运行模式。

\begin{figure}[!htbp]
\centering
\includegraphics[width=0.95\linewidth]{\detokenize{figures/diagrams/d03-architecture.pdf}}
\caption{“知微”系统四层工程架构与离线快照流转关系}
\label{fig:architecture}
\end{figure}

\noindent
\begin{minipage}[c]{0.51\linewidth}
\setlength{\parindent}{2em}\setlength{\parfillskip}{0pt plus 0.78\linewidth}
\textbf{第一层（L1）为原始字节冻结与规范数据层}。所有采集落盘的 JSONL 原始响应与评测日志存放在 \texttt{data/raw/} 下并通过 Git LFS 管理；经脱敏与类型标准化后的分实验清洗表与 21 列总表保存在 \texttt{data/processed/}，由 \texttt{SHA256SUMS.txt} 锁死字节完整性。

\textbf{第二层（L2）为离线清洗、统计推断与科研绘图层}。分布在 \texttt{scripts/} 与 \texttt{experiments/} 下的 83 个 Python 脚本（合计 17,466 行代码，源码规模分布见图~\ref{fig:loc-arch}）承担全部耗时计算，包括 AST 代码规范化、沙箱隔离执行、5,000 次按题聚类 Bootstrap 抽样以及出版级矢量图导出。

\textbf{第三层（L3）为静态构建与防漂移快照层}。构建脚本 \texttt{build\_web\_app.py} 在校验源数据与图稿 SHA-256 指纹一致后，将前端所需的统计结果打包为 \texttt{app/data/} 下仅数兆大小的匿名 JSON 快照，并为静态资源注入 12 位内容哈希后缀，彻底消除浏览器缓存旧数据的隐患。
\end{minipage}\hfill
\begin{minipage}[c]{0.465\linewidth}
\centering
\includegraphics[width=\linewidth]{\detokenize{figures/charts/c08-loc-architecture.pdf}}
\captionof{figure}{全仓库 37,534 行源码分层构成统计}
\label{fig:loc-arch}
\end{minipage}

\medskip

\textbf{第四层（L4）为零外部网络依赖的交互前端层}。启动器 \texttt{app/launch.py} 仅调用 Python 3.10+ 标准库中的 \texttt{http.server} 与 \texttt{socket} 模块即可自动完成构建检核、端口绑定与本地服务拉起。在浏览器端，主站（\texttt{index.html}）完全不加载任何第三方图表框架，直接由 \texttt{zhiwei-gate.js} 与 \texttt{zhiwei.js} 操控 DOM 与原生 SVG \texttt{viewBox}；而探索工作台（\texttt{explore.html}）所使用的 ECharts、D3 7.9.0、Vega 5.30 / Vega-Lite 5.23 与 Plotly.js 3.1.1 共 2.45\,MB 库文件均以固定版本离线内置于 \texttt{app/assets/vendor/} 目录并附许可证校验。整套系统在完全断网的环境下依然能够毫秒级响应全部交互操作。

\subsection{离线压缩载荷封装与标准库无状态服务机制}

为了让包含近九万条代码评分与十五万条问答特征的探索工作台能够在不架设后台数据库的条件下流畅筛选，我们在 L3 构建层与 L4 浏览器运行时之间设计了\textbf{“Gzip 分片静态特征包 + 预编译全局运行时脚本”}的轻量传输方案：

\begin{itemize}
  \item \textbf{冷热数据分层加载}：首屏与主叙事五屏所需的汇总统计、置信区间与精选案例体积控制在数百 KB 以内（如 \texttt{research-charts.json} 仅 93\,KB、\texttt{code-feedback-extension.json} 仅 212\,KB），页面打开瞬间即可完成解析与首帧渲染；而体量较大的逐样本匿名特征明细（如 1.85\,MB 的 \texttt{features.json.gz} 与 1.18\,MB 的 \texttt{code-features.json.gz}）仅在用户切入探索工作台打开箱线图或记录抽屉时按需异步拉取解压。
  \item \textbf{零构建工具链的纯标准库启动器}：入口脚本 \texttt{app/launch.py} 不依赖 Flask、FastAPI 或任何第三方 Web 容器。启动时，它先调用内部构建器核验 \texttt{app/dist/} 快照指纹，随后通过标准库 \texttt{http.server.ThreadingHTTPServer} 在本地环回地址 \texttt{127.0.0.1} 上挂载多线程静态文件服务，自动配置 \texttt{.wasm}、\texttt{.woff2} 与 \texttt{.svg} 的标准 MIME 类型并处理端口冲突提示，真正做到克隆即运行。
\end{itemize}
